%%%%%%%%%%  scan idx 398  |  folio 385  %%%%%%%%%%
\SGAPrintPage{385}%
it suffices, by the definition of $\bar f^{*}_{U,s}$ as left adjoint to the direct-image functor, to define an $f_{U,s*}(\underline{O}_X)$-linear homomorphism
\[
  \underline{M}_{(s),U}\longrightarrow f_{U,s*}\bigl(\underline{F}|D_+(U,s)\bigr).
\]
For this, we must give, for every open $V$ of $U$, a homomorphism:
\begin{equation}\tag{2.5.2}
  \Gamma\bigl(V,\underline{M}\bigr)_{(s)}\longrightarrow \Gamma\bigl(D_+(V,s),\underline{F}\bigr)
\end{equation}
compatible with restrictions to the opens of $U$, and take the corresponding homomorphism on the associated sheaves. Homomorphism (2.5.2) is defined as follows: let $z/s^n\in\Gamma(V,\underline{M})_{(s)}$, where, $s$ being of degree $d$, $z\in\Gamma(V,\underline{M}_{nd})=\Gamma\bigl(D_+(V,s),\underline{F}(nd)\bigr)$; we then associate to $z/s^n$ the section $z.\alpha_{nd}(s^n)^{-1}$ of $\underline{F}$ over $D_+(V,s)$, writing $\alpha_d$ for the canonical homomorphism $\Gamma(V,\underline{B}_d)\longrightarrow \Gamma\bigl(f^{-1}(V),\underline{O}_X(d)\bigr)$ (cf. EGA II 2.6.3 and 2.6.4). We easily check the necessary compatibilities, whence, by gluing, homomorphism (2.5.1).

It is clear that, in the case where $\underline{F}$ is quasi-coherent, the homomorphism thus constructed is the one in EGA II (3.3.4.1).

The homomorphism of functors $\Projz\circ F_k\to \operatorname{Id}$ on the category of $\underline{O}_X$-Modules gives, by passage to the derived categories, a homomorphism of functors on $D^+(\mathcal A)$:
\sgaStableTarget{sga6.e06.eq.2-5-2}\begin{equation}\tag{2.5.3}
  \Projz\circ\mathbb R F_k\longrightarrow \operatorname{Id}_{D^+(\mathcal A)}.
\end{equation}

\sgaStableTarget{sga6.e06.prop.2-6}\textbf{Proposition 2.6.} \uline{Let} $S$ \uline{be a quasi}-\uline{compact scheme}, $\underline{E}$ \uline{a locally free} $\underline{O}_S$\uline{-Module of finite type}, $X=\mathbf P(\underline{E})$, \uline{and} $\underline{L}^{\bullet}$ \uline{a perfect complex on} $X$. \uline{Then} \uline{there exists an integer} $k_0$ \uline{such that, for} $k\geq k_0$, \uline{the following properties hold}:
%%%%%%%%%%  scan idx 399  |  folio 386  %%%%%%%%%%
\SGAPrintPage{386}%
\[
\begin{array}{ll}
\mathrm{i)}& \mathbb R F_k(\underline{L}^{\bullet})\text{ \uline{is} }\mathcal C_0
\text{\uline{-perfect}};\\[2mm]
\mathrm{ii)}& \text{\uline{The canonical morphism} }(2.5.3):
\Projz\circ\mathbb R F_k(\underline{L}^{\bullet})\longrightarrow \underline{L}^{\bullet}
\text{ \uline{is an isomorphism in} }D^+(\mathcal A).
\end{array}
\]

The notion of a $\mathcal C_0$-perfect complex considered here is local, that is, it is perfection relative to the category fibered over $S$ defined by $\mathcal C_0$. This notion of perfection does have a meaning, since $\mathcal C_0$ satisfies the axioms of (I 4.0): only stability under kernels of epimorphisms is not immediate; it follows from 1.5.1.

Since $S$ is quasi-compact, it suffices to check the two assertions of 2.6 locally on $S$, so that we can assume $S$ affine. The invertible sheaf $\underline{O}_X(1)$ is then ample, and global resolutions on $X$ are available (II 2.2.9). There therefore exists a strictly perfect complex $\underline{L}'{}^{\bullet}$ in $D_{\mathrm{parf}}(\mathcal A)$, and an isomorphism in $D_{\mathrm{parf}}(\mathcal A)$: $\underline{L}^{\bullet}\xleftarrow{\sim}\underline{L}'{}^{\bullet}$, which allows us to assume $\underline{L}^{\bullet}$ strictly perfect. We saw in 2.4 that, for $k$ sufficiently large, the $\underline{L}^i$, all but finitely many of which are zero, are $F_k$-acyclic; we consequently have:
\[
  \mathbb R F_k(\underline{L}^{\bullet})=F_k(\underline{L}^{\bullet})
\]
for $k$ sufficiently large. By 1.3 and 1.6, we can find $k_0$ such that, for $k\geq k_0$, the preceding condition is satisfied, and the $F_k(\underline{L}^i)$ are $\mathcal C_0$-perfect $\underline{B}$-Modules for every $i$. The complex $F_k(\underline{L}^{\bullet})$ is then $\mathcal C_0$-perfect, whence i).

For $k$ chosen as above, homomorphism (2.5.3) is then the homomorphism:
%%%%%%%%%%  scan idx 400  |  folio 387  %%%%%%%%%%
\SGAPrintPage{387}%
\[
  \Projz\bigl(F_k(\underline{L}^{\bullet})\bigr)\longrightarrow \underline{L}^{\bullet},
\]
and it is an isomorphism on each component, since $\underline{L}^{\bullet}$ was assumed strictly perfect, whence ii).

\sgaStableTarget{sga6.e06.sec.2.7}2.7. On $\mathcal C$, the functor which associates to a graded $\underline{B}$-Module its $n$-th component is an exact functor with values in $\mathcal C'$; it therefore defines a functor:
\sgaStableTarget{sga6.e06.eq.2-7-1}\begin{equation}\tag{2.7.1}
  \operatorname{Comp}_n:D^+(\mathcal C)\longrightarrow D^+(\mathcal C').
\end{equation}
For every $n$, we can consider the functor $\mathcal C'\to\mathcal C$ defined by tensoring with $\underline{B}(-n)$. Homomorphism (1.2.1) defines a morphism of functors on $\mathcal C$, whence a morphism for the corresponding functors between derived categories:
\sgaStableTarget{sga6.e06.eq.2-7-2}\begin{equation}\tag{2.7.2}
  \operatorname{Comp}_n\otimes_{\underline{O}_S} \underline{B}(-n)
  \longrightarrow \operatorname{Id}_{D^+(\mathcal C)}.
\end{equation}
Let $\underline{M}^{\bullet}$ be a $\mathcal C_0$-perfect complex. We shall show that there exists a sequence of distinguished triangles:
\[
\begin{array}{ccccccccccc}
& \underline{P}_1^{\bullet} &&&& \underline{P}_2^{\bullet} &&&& \underline{P}_k^{\bullet} & \\
\swarrow && \nwarrow && \swarrow && \nwarrow && \swarrow && \nwarrow \\
\underline{N}_1^{\bullet} & \longrightarrow & \underline{M}^{\bullet}, & \qquad &
\underline{N}_2^{\bullet} & \longrightarrow & \underline{P}_1^{\bullet}, & ...., &
\underline{N}_k^{\bullet} & \longrightarrow & \underline{P}_{k-1}^{\bullet},
\end{array}
\]
where $\underline{P}_k^{\bullet}$ is acyclic, and the $\underline{N}_i^{\bullet}$ are of the form \(\underline{K}_i^{\bullet}\otimes_{\underline{O}_S}\underline{B}(d)\), the $\underline{K}_i^{\bullet}$ being perfect complexes in the ordinary sense on $S$.

Since $\underline{M}^{\bullet}$ is $\mathcal C_0$-perfect, and $S$ is quasi-compact, there exists a finite covering of $S$, and an integer $d_0$ such that, on every open of this covering, $\underline{M}^{\bullet}$ is isomorphic to a bounded complex whose terms are locally free of finite type over $\underline{B}$, say $\underline{M}_U'{}^{\bullet}$, such that $\operatorname{Comp}_d(\underline{M}_U'{}^{\bullet})=0$ for $d<d_0$.

%%%%%%%%%%  scan idx 401  |  folio 388  %%%%%%%%%%
\SGAPrintPage{388}%
We then consider the canonical morphism (2.7.2):
\[
  \underline{N}_1^{\bullet}=\operatorname{Comp}_{d_0}(\underline{M}^{\bullet})\otimes_{\underline{O}_S}\underline{B}(-d_0)
  \longrightarrow \underline{M}^{\bullet};
\]
let $\underline{P}_1^{\bullet}$ be the cone of this morphism. It is again a $\mathcal C_0$-perfect complex: indeed, for every $n$, $\operatorname{Comp}_n(\underline{M}^{\bullet})$ is a perfect complex on $S$, since this holds when $\underline{M}^{\bullet}$ is bounded with terms locally free of finite type over $\underline{B}$; consequently, the complex $\operatorname{Comp}_n(\underline{M}^{\bullet})\otimes_{\underline{O}_S}\underline{B}(-d_0)$ is $\mathcal C_0$-perfect, and the same therefore holds for the cone (I 4.10).

Moreover, on every open $U$ of the covering we can write the commutative diagram:
\[
\begin{tikzcd}[column sep=large,row sep=large]
\operatorname{Comp}_{d_0}(\underline{M}^{\bullet})\otimes_{\underline{O}_S}\underline{B}(-d_0) \arrow[r] \arrow[d,dash,"\wr"] & \underline{M}^{\bullet} \arrow[d,dash,"\wr"]\\
\operatorname{Comp}_{d_0}(\underline{M}_U'{}^{\bullet})\otimes_{\underline{O}_S}\underline{B}(-d_0) \arrow[r] & \underline{M}_U'{}^{\bullet},
\end{tikzcd}
\]
which induces an isomorphism between the cones; the lower homomorphism is an isomorphism in degree $d_0$; therefore $\underline{P}_1^{\bullet}$ is isomorphic over $U$ to a bounded complex with terms locally free of finite type over $\underline{B}$, whose components $\operatorname{Comp}_d$ are zero for $d<d_0+1$. Consequently, $\underline{P}_1^{\bullet}$ has the same properties as $\underline{M}^{\bullet}$, with $d_0$ replaced by $d_0+1$. We can therefore iterate the process, and it suffices to see that, after finitely many operations, the resulting complex $\underline{P}_k^{\bullet}$ is acyclic. Here again the question is local, so it suffices to check this for $\underline{M}_U'{}^{\bullet}$. For each term $\underline{M}_U'{}^i$, the morphisms considered above are those defining the filtration of Proposition 1.7; this filtration is finite, so, starting from $\underline{M}_U'{}^{\bullet}$, we reach the zero complex after finitely many operations, that is, the corresponding complex $\underline{P}_k^{\bullet}$ is acyclic.

%%%%%%%%%%  scan idx 402  |  folio 389  %%%%%%%%%%
\SGAPrintPage{389}%
\sgaStableTarget{sga6.e06.sec.2.8}2.8.
Let $\underline{L}^{\bullet}$ be a perfect complex on $X$. There exists an integer $k$ such that $\underline{M}^{\bullet}=\mathbb R F_k(\underline{L}^{\bullet})$ is a $\mathcal C_0$-perfect complex, and $\underline{L}^{\bullet}$ is isomorphic to $\Projz\bigl(\mathbb R F_k(\underline{L}^{\bullet})\bigr)$. Applying 2.7 to the complex $\underline{M}^{\bullet}$, the sequence of distinguished triangles constructed gives, on applying the functor $\Projz$, a sequence of distinguished triangles in $D_{\mathrm{parf}}(\mathcal A)$ (it is clear that $\Projz$ transforms a $\mathcal C_0$-perfect complex into a perfect complex). Moreover, we have:
\[
  \Projz\bigl(\underline{K}^{\bullet}\otimes_{\underline{O}_S}\underline{B}(d)\bigr)
  = f^*(\underline{K}^{\bullet})(d)
\]
(note that $f^*=\mathbb L f^*$ by the flatness of $f$) for $\underline{K}^{\bullet}$ in $D^+(\mathcal C')$.

We therefore obtain in $K^{\bullet}(X)$ the equality:
\[
  \cl(\underline{L}^{\bullet})
  =\sum_i\cl\bigl(f^*(\underline{K}_i^{\bullet})\bigr).\xi^{d_i}
  =\sum_i f^*\bigl(\cl(\underline{K}_i^{\bullet})\bigr).\xi^{d_i},
\]
where the $\underline{K}_i^{\bullet}$ are the perfect complexes on $S$ defined in 2.7. It follows that $K^{\bullet}(X)$ is generated as a $K^{\bullet}(S)$-module by the powers of $\xi$.

To see that $1,\xi,\ldots,\xi^{r+1}$ are linearly dependent, we consider the homomorphism: $K^{\bullet}_{\naif}(X)\to K^{\bullet}_{\mathrm{parf}}(X)$ which associates to the class of a locally free sheaf of finite type $\underline{L}$ the class of the complex equal to $\underline{L}$ in degree $0$, and zero in the other degrees. Equation (1.9.1) then gives a relation of dependence in $K^{\bullet}_{\mathrm{parf}}(X)$.

To see that $1,\xi,\ldots,\xi^r$ are linearly independent, we can use the $K^{\bullet}(S)$-linear morphism $f_*:K^{\bullet}(X)\to K^{\bullet}(S)$ (IV 2.11.1). We have:
\[
  f_*(1)=1,
  \quad\text{since } f_*(\underline{O}_X)=\underline{O}_S,\text{ and }R^i(f_*(\underline{O}_X))=0\text{ for }i\geq1;
\]
\[
  f_*(\xi^n)=0\quad\text{for }-r\leq n<0,
  \quad\text{since }R^i\bigl(f_*(\underline{O}_X(n))\bigr)=0\text{ for }i\geq0.
\]
Given a relation of the form:
\[
  f^*(x_0)+f^*(x_1)\xi+\ldots+f^*(x_r)\xi^r=0,
\]
we multiply it by $\xi^{-r}$, and apply $f_*$; we then find $x_r=0$; we
